\subsection{\texorpdfstring{环 \(A^{(d)}\) 的性质}{环 A(d) 的性质}}

设 \(\mathbf{A} = \bigoplus_{i\in \mathbf{Z}}\mathbf{A}_i\) 是一个分次环，\(\mathbf{M} = \bigoplus_{i\in \mathbf{Z}}\mathbf{M}_i\) 是一个分次 A-模；对每一对有序整数 \((d, k)\) （满足 \(d \geqslant 1\) 与 \(0 \leqslant k \leqslant d - 1\) ），令

\[
\mathbf {A} ^ {(d)} = \bigoplus_ {i \in \mathbf {Z}} \mathbf {A} _ {i d}, \quad \mathbf {M} ^ {(d, k)} = \bigoplus_ {i \in \mathbf {Z}} \mathbf {M} _ {i d + k}.
\]

显然 \(\mathbf{A}^{(d)}\) 是 A 的一个分次子环，\(\mathbf{M}^{(d,k)}\) 是一个分次 \(\mathbf{A}^{(d)}\) -模；此外，若 N 是 M 的一个分次子模，则 \(\mathbf{N}^{(d,k)}\) 是分次 \(\mathbf{A}^{(d)}\) -模 \(\mathbf{M}^{(d,k)}\) 的一个子模。我们将用 \(\mathbf{M}^{(d)}\) 代替 \(\mathbf{M}^{(d,0)}\) 来写；对每个 \(d \geqslant 1\) ，M 是诸 \(\mathbf{A}^{(d)}\) -模 \(\mathbf{M}^{(d,k)}\) （ \(0 \leqslant k \leqslant d - 1\) ）的直和。

命题 2. 设 \(\mathbf{A} = \bigoplus_{i\in \mathbf{Z}}\mathbf{A}_i\) 是一个正次数的分次交换环，\(\mathbf{M} = \bigoplus_{i\in \mathbf{Z}}\mathbf{M}_i\) 是一个分次 A-模。假设 \(\mathbf{A}\) 是一个有限生成 \(\mathbf{A}_0\) -代数，而 \(\mathbf{M}\) 是一个有限生成 A-模。则对每一对有序整数 \((d,k)\) （满足 \(d\geqslant 1,0\leqslant k\leqslant d - 1\) ）：

\begin{enumerate}
\def\labelenumi{(\roman{enumi})}
\item
  \(\mathbf{A}^{(d)}\) 是一个有限生成 \(\mathbf{A}_0\) -模。
\item
  \(\mathbf{M}^{(d,k)}\) 是一个有限生成 \(\mathbf{A}^{(d)}\) -模。
\end{enumerate}

我们来证明 A 是一个有限生成 \(\mathbf{A}^{(d)}\) -模。设 \((a_{i})_{1\leqslant i\leqslant s}\) 是 \(\mathbf{A}_0\) -代数 A 的由齐次元组成的一个生成系。A 中形如 \(a_1^{\alpha_1}a_2^{\alpha_2}\ldots a_s^{\alpha_s}\) 且满足 \(0\leqslant \alpha_{i}\leqslant d\) （对 \(1\leqslant i\leqslant s\) ）的元素（数目有限）组成 \(\mathbf{A}^{(d)}\) -模 A 的一个生成系；对每一组整数 \(n_i\geqslant 0\) （ \(1\leqslant i\leqslant s\) ），存在正整数 \(q_{i}, r_{i}\) ，使得 \(n_i = q_id + r_i\) ，其中 \(r_i < d\) （ \(1\leqslant i\leqslant s\) ）；于是

\[
a _ {1} ^ {n _ {1}} a _ {2} ^ {n _ {2}} \dots a _ {s} ^ {n _ {s}} = (a _ {1} ^ {q _ {1}} \dots a _ {s} ^ {q _ {s}}) ^ {d} (a _ {1} ^ {r _ {1}} \dots a _ {s} ^ {r _ {s}})
\]

这就证明了我们的断言，因为每个齐次元 \(x \in A\) 都满足 \(x^{d} \in \mathrm{A}^{(d)}\) 。于是，若 M 是一个有限生成 A-模，则它也是一个有限生成 \(\mathrm{A}^{(d)}\) -模；由于 M 是诸 \(\mathrm{M}^{(d,k)}\) （ \(0 \leqslant k \leqslant d - 1\) ）的直和，诸 \(\mathrm{M}^{(d,k)}\) 中的每一个都是有限生成 \(\mathrm{A}^{(d)}\) -模，这就证明了 (ii)。

把上述结果应用于分次 A-模 \(m = \bigoplus_{i \geqslant 1} A_i\) ，依第 2 节命题 1 的推论它是有限生成的；可以看出 \(m^{(d)}\) 是一个有限生成 \(A^{(d)}\) -模；因此（第 2 节，命题 1 的推论）\(A^{(d)}\) 是一个有限生成 \(A_0\) -代数。

引理 1. 设 A 是一个满足 \(\mathbf{A} = \mathbf{A}_0[\mathbf{A}_1]\) 的分次交换环，M 是一个分次 A-模，\((y_{\lambda})_{\lambda \in \mathbf{L}}\) 是 M 的一个齐次生成系，使得 \(\deg(y_{\lambda}) \leqslant n_0\) 对所有 \(\lambda \in \mathbf{L}\) 成立。则对所有 \(n \geqslant n_0\) 及所有 \(k \geqslant 0\) ，有 \(\mathbf{M}_{n+k} = \mathbf{A}_k \cdot \mathbf{M}_n\) 。

设 \(n \geqslant n_0\) ，\(x \in \mathbf{M}_{n+1}\) 。由于诸 \(y_\lambda\) 生成 \(\mathbf{M}\) ，存在一个有限支撑族 \((a_\lambda)_{\lambda \in \mathbf{L}}\) ，其元素属于 \(\mathbf{A}\) ，使得 \(x = \sum_{\lambda} a_\lambda y_\lambda\) ；进一步可以假设每个 \(a_\lambda\) 都是次数为 \(n + 1 - \deg(y_\lambda)\) 的齐次元（必要时用该次数的齐次分量代替它）。由于

\(A = A_{0}[A_{1}]\) 且 \(\deg(a_{\lambda}) > 0\) ，每个 \(a_{\lambda}\) 都是形如 \(bb'\) 的元素之和，其中 \(b \in A_{1}\) ，\(b' \in A\) ，从而 \(x \in A_{1}M_{n}\) 。于是 \(M_{n+1} = A_{1}M_{n}\) ，从而对 k 作归纳得 \(M_{n+k} = A_{k}M_{n}\) 。

引理 2. 设 A 是一个满足 \(\mathbf{A} = \mathbf{A}_0[\mathbf{A}_1]\) 的分次交换环，\(\mathbf{S} = \bigoplus_{i\geqslant 0}\mathbf{S}_i\) 是一个正次数的分次交换 A-代数，且是有限生成 A-模。则存在一个整数 \(n_0\geqslant 0\) ，使得：

\begin{enumerate}
\def\labelenumi{(\roman{enumi})}
\item
  对 \(n \geqslant n_0\) 与 \(k \geqslant 0\) ，\(\mathrm{S}_{n+k} = \mathrm{S}_k \cdot \mathrm{S}_n\) 。
\item
  对 \(d \geqslant n_0\) ，\(\mathbf{S}^{(d)} = \mathbf{S}_0[\mathbf{S}_d]\) 。
\end{enumerate}

依引理 1，存在一个整数 \(n_0 \geqslant 0\) ，使得对 \(n \geqslant n_0\) 与 \(k \geqslant 0\) 有 \(S_{n+k} = A_kS_n\) ，更有 \(S_{n+k} = S_kS_n\) ，这就建立了 (i)。于是，对 \(d \geqslant n_0\) 与 \(m > 0\) ，\(S_{md} = (S_d)^m\) （应用 (i) 并对 \(m\) 作归纳即得）；这就建立了 (ii)。

命题 3. 设 \(\mathbf{R} = \bigoplus_{i\geqslant 0}\mathbf{R}_i\) 是一个正次数的分次交换环，且是有限生成 \(\mathbf{R}_0\) -代数。则存在一个整数 \(e\geqslant 1\) ，使得 \(\mathbf{R}^{(me)} = \mathbf{R}_0[\mathbf{R}_{me}]\) 对所有 \(m\geqslant 1\) 成立。

设 \((x_{j})_{1\leqslant j\leqslant s}\) 是 R \(_{0}\) -代数 R 的一个由次数 \(\geqslant 1\) 的齐次元组成的生成系。设 \(h_{j} = \deg(x_{j})\) ，设 \(q\) 是诸 \(h_{j}\) 的一个公倍数，并记 \(q_{j} = q/h_{j}\) （ \(1 \leqslant j \leqslant s\) ）；则诸元素 \(x_{j}^{q_{j}}\) 的次数都等于 \(q\) 。设 B 是 R 的一个子 R \(_{0}\) -代数，由诸 \(x_{j}^{q_{j}}\) 生成；它是 R 的一个分次子代数，且 B \(_{i} = 0\) （若 \(i\) 不是 \(q\) 的倍数）。设 A（相应地 S）是这样的分次环：其底层环为 B（相应地 R \(^{(q)}\) ），其分次为 A \(_{i} =\) B \(_{iq}\) （相应地 S \(_{i} =\) R \(_{iq}\) ）。依 B 的定义，有 A = A \(_{0}\) {[}A \(_{1}\) {]} 。考虑 R 中形如 \(x_{1}^{\alpha_{1}}x_{2}^{\alpha_{2}}\ldots x_{s}^{\alpha_{s}}\) 的（有限多个）元素，其中 \(0 \leqslant \alpha_{j} \leqslant q_{j}\) 且 \(\alpha_{1}h_{1} + \cdots + \alpha_{s}h_{s} \equiv 0 \pmod{q}\) ；我们来证明它们生成 B-模 R \(^{(q)}\) 。只需证明 R \(^{(q)}\) 的每个形如 \(x_{1}^{n_{1}}x_{2}^{n_{2}}\ldots x_{s}^{n_{s}}\) 的元素都是上述元素的一个 B-线性组合。现在，存在正整数 \(k_{j}, r_{j}\) ，使得 \(n_{j} = k_{j}q_{j} + r_{j}\) ，其中 \(r_{j} < q_{j}\) （ \(1 \leqslant j \leqslant s\) ）；于是

\[
x _ {1} ^ {n _ {1}} x _ {2} ^ {n _ {2}} \dots x _ {s} ^ {n _ {s}} = (x _ {1} ^ {q _ {1}}) ^ {k _ {1}} \dots (x _ {s} ^ {q _ {s}}) ^ {k _ {s}}. (x _ {1} ^ {r _ {1}} \dots x _ {s} ^ {r _ {s}})
\]

又由假设 \(\sum_{j=1}^{s} n_j h_j \equiv 0 \pmod{q}\) ，故 \(\sum_{j=1}^{s} r_j h_j \equiv 0 \pmod{q}\) ；由于依定义诸 \(x_j^{q_j}\) 属于 B，这就证明了我们的断言。由于 S 是一个有限生成 A-模，可以应用引理 2：存在 \(n_0\) ，使得对 \(d \geqslant n_0\) 有 \(\mathrm{S}^{(d)} = \mathrm{S}_0[\mathrm{S}_d]\) ，从而 \(\mathrm{R}^{(qd)} = \mathrm{R}_0[\mathrm{R}_{qd}]\) 对 \(d \geqslant n_0\) 成立。取 \(e = q n_0\) 即得命题。

